\clearpage
\section{Validation, artifacts, and reproduction}\label{sec:validation}

\subsection{What has actually been checked}

\begin{center}
\begin{tabularx}{\linewidth}{@{}L{.28\linewidth}YL{.20\linewidth}@{}}
\toprule
Item & Evidence & Status\\
\midrule
Finite trace theorem & General proof in Section~\ref{sec:trace}; nonnormal and repeated-root examples & Proved mathematically\\
Bordered integral and source adapter & General proof in Section~\ref{sec:hua}; exact preservation of the majorant constant & Proved mathematically\\
Finite interpolation and trace transport & Two standalone Lean files with pinned dependencies; compiler logs and printed axiom sets & Kernel-checked support\\
Full spectral trace certificate & Written proof supplied; formal support takes the certificate as an explicit premise & Not kernel-checked here\\
Real-place source replacement & Source interfaces and call sites inspected; exact algebra and geometry regression checks & Not implemented in GRH\\
All public GRH endpoints & Source statements inspected; no new full-project build or independent kernel replay & Not reverified here\\
\bottomrule
\end{tabularx}
\end{center}

The distinction between the third and fourth rows is essential. The checked
theorem \code{trace_projector_eq_of_polynomial_certificates} explicitly assumes
the two identities \(\tr(Ae)=\tr(Ap(T))\) and
\(\tr(A'e')=\tr(A'p(T'))\). It proves the finite moment transport between them.
It does not silently derive those identities from the line relation.

\subsection{Standalone Lean support}

\code{lean/SpectralTraceCore.lean} proves five statements: a zero-constant
finite interpolant with the degree bound; positive-moment transport for linear
functionals; its weighted-trace specialization; logarithmic rank extraction
from the scalar determinant-norm formula; and the final conditional trace transfer.

\code{lean/SpectralTraceSpectrum.lean} proves that a nonzero determinant excludes
zero from the characteristic roots, bounds the number of distinct roots by the
dimension, and constructs a common interpolant of degree at most
\(\dim V+\dim V'\). Over \(\C\), these characteristic roots are the eigenvalues
used in the written theorem.

Both files compiled with Lean 4.35.0 and Mathlib commit
\code{56a005855ea17e0b0958f73ec801ae140b0a1679}, with exit status zero and no
warnings. The six printed endpoints depend only on
\code{propext}, \code{Classical.choice}, and \code{Quot.sound}.
There are no supplied \code{sorry} proofs or added axioms.
The accompanying manifest records file hashes and the exact verification scope.

The archive contains the pinned Lake configuration and a portable README.
The supplied compiler logs come from checking these files against an existing
cache at the pinned Mathlib revision, with a separate overlay for missing compiled
Mathlib modules. They are not logs of a complete GRH build.

\Needspace{12\baselineskip}
\subsection{Independent exact-arithmetic checks}

The checks complement the general proofs; finite examples do not prove the
universal statements.

\begin{itemize}
\item \code{check_trace_transfer.py}: 107 assertions over 11 examples.
They include a nonnormal Jordan operator, repeated roots, nonscalar and
nonnormal operators in the common commutant, simultaneous interpolation,
the failure when adjoint commutation is dropped, and the failure of a
single positive moment. Rational matrix identities are checked exactly.
One additional logarithmic evaluation is explicitly checked numerically at
80-digit precision and is labelled as such in the output.

\item \code{check_bordered_hua.py}: 36 positive rational-complex matrix examples
with \(1\leq k\leq3\), \(0\leq i\leq3\), and three examples per dimension pair.
Each checks eight determinant/congruence identities, including the adapter to
the existing Hua bound. Nine exact real-Jacobian checks reach real dimension 18.
Symbolic calculations verify the exponent cancellations and the final source
constant; a scalar radial integral is exactly \(\pi/28\).

\item \code{check_bordered_geometry.py}: 248 exact assertions in 11 rational
\(J\)-unitary examples shaped like the source geometry, including \(i=0\),
\(k=7\), both level branches, nontrivial subframe determinants, and Bessel
unipotent translations. These check the canonical dual frame, direct determinant
reduction, energy identity, and unchanged majorant factors.

\item \code{audit_sources.py}: verifies the pinned repository identity,
803-module/476,849-line inventory, static import closure, proposal-bundle count,
target-file hashes, and affected source regions. This is a reproducible source
inventory, not a proof-term dependency analysis.
\end{itemize}

All three mathematical check scripts run without network access and require
Python with SymPy 1.14.0. The source audit uses Python's standard library and
a local checkout of the pinned repository. Commands and outputs are included
in the archive README and \code{results/}.

\subsection{Archive contents}

\begin{center}
\begin{tabularx}{\linewidth}{@{}L{.32\linewidth}Y@{}}
\toprule
Artifact & Purpose\\
\midrule
\code{paper.tex}, \code{sections/} & Complete editable report, with all proofs and references\\
\code{paper.pdf} & Rendered report\\
\code{lean/} & Checked support files, pinned dependencies, reproduction instructions\\
\code{checks/} & Executable algebra, geometry, and source-audit checks\\
\code{results/} & JSON results, Lean logs, provenance, and file hashes\\
\code{audit/} & Proposal comparison matrix and proposal-text manifest\\
\code{README.md}, \code{Makefile} & Entry points for readers and reproduction\\
\bottomrule
\end{tabularx}
\end{center}

The main implementation recommendation is therefore specific: retain the
unnormalized negative dual frame, integrate its bordered Gram determinant using
the existing Hua theorem, and preserve the public majorant statement exactly.
The finite trace theorem is an independent second improvement whose mathematical
contract is sharper, with its remaining formalization cost explicitly identified.
